% ======================================================================
\section{Introduction}\label{sec:intro}
% ======================================================================

At high energy the wave function of a hadron contains a large number of soft gluons, and in a sufficiently dense system their density saturates~\cite{Gribov:1983ivg,Mueller:1985wy}. The Color Glass Condensate (CGC) is the effective theory of this regime~\cite{McLerran:1993ni,McLerran:1993ka,Iancu:2003xm,Weigert:2005us,Gelis:2010nm,Kovchegov:2012mbw}: the fast (valence) degrees of freedom of a hadron act as static colour sources, characterised by a colour charge density $\rho^a(\v x)$, for the soft gluon field, and the dependence of observables on the rapidity interval over which the soft modes are integrated out is governed by the BFKL equation~\cite{Kuraev:1977fs,Balitsky:1978ic} in the dilute regime and by the non-linear Balitsky--Kovchegov (BK) and JIMWLK equations~\cite{Balitsky:1995ub,Kovchegov:1999yj,Kovchegov:1999ua,Jalilian-Marian:1997qno,Jalilian-Marian:1997jx,Jalilian-Marian:1997ubg,Kovner:2000pt,Weigert:2000gi,Iancu:2000hn,Iancu:2001ad,Ferreiro:2001qy} in the dense regime.

The simplest observable that is sensitive both to the projectile and to the target wave functions is the single inclusive gluon spectrum in the collision of a dilute projectile with a dense target. At leading order (LO), for a gluon produced at mid rapidity, i.e.\ far in rapidity from both the projectile and the target, the spectrum is given by the formula of Kovchegov and Tuchin~\cite{Kovchegov:2001sc,Kovchegov:1998bi}: the gluon is emitted from the valence charges of the projectile in the Weizs\"acker--Williams field, and the emission amplitude, together with the colour rotation of the emitting charge, scatters eikonally on the target. The cross section is a product of the projectile colour charge correlator and of a target operator built from two adjoint Wilson lines; in the limit of a dilute target it reduces to the $k_T$-factorised form with the Lipatov vertex~\cite{Gribov:1983ivg,Kharzeev:2003wz,Blaizot:2004wu}. The leading logarithmic (LL) rapidity evolution of this result was obtained by Kovner and Lublinsky~\cite{Kovner:2006wr} (see also~\cite{Kovner:2006ge}) and, within the Schwinger--Keldysh formulation, by Gelis, Lappi and Venugopalan~\cite{Gelis:2008rw}: the projectile is to be evolved from its own rapidity down to the rapidity of the produced gluon, and the target from the rapidity of the produced gluon to its own rapidity. For a dilute projectile the first evolution is linear (BFKL) and acts on the projectile charge correlator, while the second is the JIMWLK evolution of the target Wilson lines.

To make quantitative use of this result, the next-to-leading order (NLO) corrections, i.e.\ the terms of relative order $\alpha_s$ that are not enhanced by a rapidity logarithm, are required. The ingredients of such a calculation have become available over the last decade. The evolution equations themselves are known at NLO: the NLO BK equation~\cite{Balitsky:2007feb} and the NLO JIMWLK Hamiltonian~\cite{Balitsky:2013fea,Kovner:2013ona,Kovner:2014lca,Grabovsky:2013mba}, together with the running-coupling corrections to the LO kernels~\cite{Balitsky:2006wa,Kovchegov:2006vj,Kovchegov:2006wf}. In particular, Lublinsky and Mulian~\cite{Lublinsky:2016meo} computed the soft gluon light-cone wave function of a dilute projectile to order $g^3$ and derived the NLO JIMWLK Hamiltonian from it. For particle production, the NLO corrections to single inclusive hadron production at forward rapidity were computed in the hybrid formalism~\cite{Dumitru:2005gt,Altinoluk:2011qy} by Chirilli, Xiao and Yuan~\cite{Chirilli:2011km,Chirilli:2012jd}, and the structure of the rapidity subtraction and of the kinematic constraints in this framework was analysed in~\cite{Stasto:2013cha,Altinoluk:2014eka,Ducloue:2016shw}. Running-coupling corrections to inclusive gluon production at mid rapidity were obtained by Horowitz and Kovchegov~\cite{Horowitz:2010yg}, and the relation between $\beta_0$-dependent terms of NLO calculations and the running of the coupling was clarified in~\cite{Kovner:2023vsy}. The general structure of the NLO correction to inclusive spectra in the CGC was discussed by Gelis and Laidet~\cite{Gelis:2012ri}. Finally, corrections of higher order in the projectile colour charge (``first saturation corrections'' in the projectile) were computed in the classical approximation in~\cite{Chirilli:2015tea,Li:2021zmf,Li:2021ntm}, and their contribution to single inclusive production, odd in $\v k$, was analysed in~\cite{McLerran:2016snu,Kovchegov:2018jun}. A complete NLO calculation of the single inclusive gluon spectrum at mid rapidity has, to our knowledge, not been presented.

In a previous paper~\cite{Radhakrishnan:2026soft} the soft gluon wave function of a dilute projectile and the associated unitary evolution (dressing) operator $\Omega$ were constructed up to order $g^2$, in pure Yang--Mills theory, within light-cone perturbation theory (LCPT)~\cite{Lepage:1980fj,Brodsky:1997de}. The NLO cross section requires in addition the coefficient of $\Omega$ that creates one gluon at order $g^3$; its ingredients, the $O(g^3)$ one-gluon components of the light-cone wave function, were computed by Lublinsky and Mulian~\cite{Lublinsky:2016meo}, and we recompute them here with the longitudinal cut-offs and the operator ordering required by the cross section. The operator $\Omega$ was written as a normal-ordered series in soft gluon creation and annihilation operators with coefficients that are functionals of the non-commuting valence charge density $\rho$, and the coefficients were fixed by matching to LCPT and by unitarity. In the present paper we use $\Omega$ to compute the single inclusive gluon spectrum at mid rapidity at NLO. The observable is the expectation value of the number operator of dressed gluons in the scattered, dressed projectile state~\cite{Kovner:2006wr}, so that the full calculation reduces to an algebraic evaluation of products of the coefficients of $\Omega$ and of Wilson lines, followed by longitudinal and transverse integrations.

Our main results are the following.
\begin{itemize}
\item The $O(g^4)$ cross section naturally decomposes into three groups of terms: interference of the LO amplitude with the $O(g^3)$ one-gluon amplitude (Group~I); the squares of the $O(g^2)$ two-gluon amplitudes, with one of the gluons integrated over (Group~II); and the interference terms between different two-gluon amplitudes (Group~III). We list all products of coefficients that contribute, with their longitudinal integration ranges and their order in $\rho$ (Section~\ref{sec:observable}).
\item All terms enhanced by $\log(\vee/\Lambda)$, where $\Lambda<p^+<\vee$ is the longitudinal momentum range of the soft modes, are the LL evolution of the LO result. The logarithm splits at the longitudinal momentum $k^+$ of the produced gluon into $\log(\vee/k^+)+\log(k^+/\Lambda)$; the coefficient of $\log(\vee/k^+)$ is exactly the BFKL evolution of the projectile charge correlator inserted into the LO formula, and the coefficient of $\log(k^+/\Lambda)$ is exactly the JIMWLK Hamiltonian acting on the target operator of the LO formula. We give both terms for an arbitrary (ordered) projectile correlator and in colour-singlet form. The contributions quartic in $\rho$, all of which carry $\log(\vee/\Lambda)$, cancel identically for commuting charges, and the totally symmetrised cubic terms cancel as well; with non-commuting charges, their reordering produces the two-$\rho$ terms that complete the evolution (Section~\ref{sec:rapidity}).
\item The remaining terms are a collinear term, associated with the fragmentation of the produced gluon and governed by the gluon--gluon splitting function, a running-coupling term, and finite terms. We give an account of all poles $1/\epsilon$: the ultraviolet poles of the virtual corrections to the produced gluon cancel among themselves, the coupling is renormalised by the normalisation of the dressed produced gluon, and the collinear poles combine into the full $\MS$ splitting function $\mathcal P_{gg}$ and are absorbed into the fragmentation function. The running-coupling term is the $\frac{\beta_0}{2}\delta(1-\xi)$ part of the collinear term; it sets the scale of the LO coupling (Section~\ref{sec:dglap_rc}).
\item We give the finite terms for an arbitrary ordered charge correlator and, averaged over the projectile, in colour-singlet form. The double logarithms $\log^2(k^+/\Lambda)$ of individual rows cancel in the sum. The soft remainders of the plus prescriptions cancel in pairs of a virtual and a real term with the same colour structure and kernel, without any condition on the target at large distance; in particular we show which terms cancel the non-integrable logarithm of one of the virtual rows. After the average over the projectile the terms cubic in $\rho$ do not vanish: their symmetrised part, odd in $\v k$, is proportional to the $d^{abc}$ (odderon) three-point function of the projectile, kept general, and the commutators of the ordered charges leave contact terms, even in $\v k$, which are present for any projectile; both are given row by row (Section~\ref{sec:finite}).
\end{itemize}
The result can be written schematically as
\begin{align}\label{eq:schematic}
\frac{d^3N}{d^3k} = \frac{d^3N_{\rm LO}}{d^3k}
+ \log\frac{\vee}{k^+}\,\frac{d^3N_{\rm BFKL}}{d^3k}
+ \log\frac{k^+}{\Lambda}\,\frac{d^3N_{\rm JIMWLK}}{d^3k}
+ \frac{d^3N_{\rm DGLAP}}{d^3k}
+ \frac{d^3N_{\rm RC}}{d^3k}
+ \frac{d^3N_{\rm finite}}{d^3k}\,,
\end{align}
where every term on the right-hand side except the first is of relative order $\alpha_s$.

The paper is organised as follows. In Section~\ref{sec:setup} we summarise the setup: light-cone conventions, the Hamiltonian of the soft modes in the eikonal approximation, LCPT with operator-valued matrix elements, and the expansion of $\Omega$ with its coefficients in coordinate space. In Section~\ref{sec:observable} we define the observable, rederive the LO cross section and organise the NLO contributions. Section~\ref{sec:rapidity} is devoted to the rapidity logarithms and their identification with the BFKL and JIMWLK evolution. The collinear and running-coupling terms are discussed in Section~\ref{sec:dglap_rc}, the finite terms in Section~\ref{sec:finite}, and limits and comparisons with the literature in Section~\ref{sec:checks}. We conclude in Section~\ref{sec:conclusions}. The appendices collect the Hamiltonian and its matrix elements (Appendix~\ref{app:hamiltonian}), the light-cone wave functions from which the coefficients of $\Omega$ are obtained (Appendix~\ref{app:wavefunctions}), the unitarity relations among the coefficients (Appendix~\ref{app:unitarity}), the Fourier transforms (Appendix~\ref{app:fourier}), the prescriptions and distributions used in the longitudinal and transverse integrals (Appendix~\ref{app:distributions}) and the complete list of the finite terms for an arbitrary ordered charge correlator (Appendix~\ref{app:finite}).
